\section{Vibe Coding Bad Behavior Across Toolchains}

Vibe coding does not invent new bad behavior. It accelerates old bad behavior and makes it easier to hide inside a plausible diff. A human can misuse \texttt{unsafe}, paste string-built SQL, disable a security job, or force-push a branch. An agent can do the same faster, across more files, with permission to run tools, rewrite tests, and produce confident explanations. The review problem is therefore not whether the author was human or synthetic. The problem is whether the repository can reject non-defensible behavior and require evidence strong enough for another reviewer or agent to replay.

The language bad-behavior detector family is the narrow version of that idea. It does not try to score taste. It identifies constructs whose safety depends on missing proof: unsound Rust, injectable or destructive SQL, TypeScript boundary lies, unsafe Python execution paths, privileged Docker builds, weak CI trust boundaries, destructive Git automation, and unsafe release automation. Those rows belong in the main paper because they make the standard concrete. A repository cannot be agent-ready if the common toolchain escape hatches are treated as harmless style.

\subsection{Rust}

Rust improves the default safety baseline, but its escape hatches are deliberately sharp. Undocumented \texttt{unsafe}, public \texttt{unsafe fn} without caller obligations, raw ownership reconstruction, unchecked indexing, fabricated UTF-8, broad lint suppression, and dynamic shell execution all move correctness out of the compiler and into a proof obligation. In vibe coding, the recurring failure is to silence the borrow checker or type checker until the program compiles, then treat compilation as evidence. That is not defensible. The reviewable form is local: every unsafe assumption needs a nearby \texttt{SAFETY:} argument or equivalent contract, tests must exercise the invalid case where possible, and shell or FFI boundaries need security proof rather than a verbal explanation.

\subsection{SQL and PostgreSQL}

SQL failures become durable quickly because the database remembers them. String-built queries turn input into code. Destructive migrations can remove data faster than reviewers can reconstruct intent. Missing constraints outsource truth to every caller forever. PostgreSQL-specific behavior also matters: null semantics, locking, migration order, isolation levels, indexes, and dialect details are not generic SQL trivia. A generated migration that works on an empty local database is not merge evidence. Jankurai expects parameter binding or identifier allowlists, database constraints for durable invariants, migration receipts with rollback or forward-repair paths, row-count and lock evidence for large tables, and tests for tenant isolation, nulls, duplicates, and concurrent writes.

\subsection{TypeScript}

TypeScript makes product surfaces easier to review only when the type boundary is honest. \texttt{as any}, double casts, disabled strictness, unchecked \texttt{JSON.parse}, raw \texttt{response.json()}, raw HTML sinks, shell execution, and raw SQL calls all create a false typed surface. The code looks typed while the runtime boundary remains unproved. That is especially risky in agent-authored UI and API glue because the model can satisfy the compiler by moving uncertainty into casts. The repair is not more type theater. The repair is a runtime decoder or generated contract at the boundary, strict compiler evidence, negative tests for malformed input, and rendered or API receipts showing the changed behavior under realistic failure cases.

\subsection{Python}

Python is not bad by identity, but dynamic execution paths are hard to defend in a repository that needs deterministic merge evidence. \texttt{eval}, \texttt{exec}, unsafe object deserialization, \texttt{shell=True}, string-built database access, TLS verification bypasses, and broad exception swallowing all erase the boundary that reviewers need. In this workspace's reference profile, Python is additionally constrained because it must not own product truth, authorization, production database writes, proof lanes, or repo tooling except under a dated advanced-ML/data exception. The reviewable form is a boxed service boundary, typed inputs and outputs, safe loaders, argv-array subprocess calls, parameterized database access, explicit exception ownership, and proof that the Python surface cannot silently become durable truth.

\subsection{Docker}

Containers are often treated as disposable build packaging, but Docker files are authority maps. \texttt{--privileged}, Docker socket mounts, host namespaces, broad capabilities, unconfined security profiles, mutable tags, public database ports, secret-bearing build layers, and remote install pipes all weaken the isolation reviewers think they have. In vibe-coded infrastructure, these changes often appear as one-line fixes for a failing build. That is not an acceptable proof. A defensible container change pins images by digest or stable version, verifies remote downloads, keeps secrets out of build context and layers, runs as least privilege, limits published ports, and leaves a security or image-scan receipt.

\subsection{CI}

CI is the repository's execution boundary. Bad CI behavior is therefore not merely operational debt; it can turn a pull request into privileged code execution. The dangerous patterns are well known: \texttt{pull\_request\_target} combined with untrusted checkout, broad \texttt{write-all} permissions, mutable action references, secret printing, artifacts or caches that include credentials, privileged Docker on untrusted runners, and security scans configured to continue on error. An agent that can edit CI can also edit the proof system unless the repository separates trusted and untrusted contexts. Jankurai expects minimum permissions, pinned actions, blocking security jobs, workflow owner review, artifact provenance, CI diff receipts, and explicit trust separation before a CI change can support a merge.

\subsection{Git and Git Tooling}

Git is safe when history and working-tree state are explicit. It is dangerous when automation hides state or mutates too broadly. Hard resets, destructive cleans, stash-based workflows, force pushes, broad \texttt{git add .}, \texttt{--no-verify}, destructive ref edits, and credential-bearing remotes all defeat review locality. Hooks, MCP servers, skills, terminal agents, and IDE agents add the same risk through another door: they can mutate repository state while looking like helper infrastructure. A reviewable Git tool names exact paths, records status before staging, avoids hidden state, refuses verification bypass, scopes destructive operations to explicit temporary paths, and leaves tool identity plus changed-path receipts.

Across these families, the pattern is stable. A bad behavior is non-defensible when it moves trust from a machine-checkable boundary into an unrecorded assumption. A Jankurai repair makes the assumption local, typed, constrained, tested, and receipted. That is why the rule IDs matter: \texttt{HLT-029} through \texttt{HLT-037} are not language rankings. They are merge-evidence obligations for the places where vibe coding most often converts speed into hidden risk.
